% ===== latex-recomposed:document =====
% flm: style file flmrecomposed.e638e2e2.sty
(see \flmRef{ref-v1}) (see \flmRef{ref-v1}) (see \flmRef{ref-v1})%%
% ===== latex-recomposed:style flmrecomposed.e638e2e2.sty =====
\NeedsTeXFormat{LaTeX2e}
\ProvidesPackage{flmrecomposed.e638e2e2}
% --- packages ---
% flm.latex.pkg.zref-clever
\RequirePackage{zref-clever}
% --- commands and environments the library defines for the document ---
% flm.latex.def.refs
% \flmRef[<options>]{<label>}: a reference that gives the name of its target
% ("Section 2"), set by \zcref of the zref-clever package.
%
% The options are the ones the document wrote, in FLM's own vocabulary. This
% definition translates each of them by its key -- a value written after the
% key is ignored, as in FLM -- and drops every key it has no translation for:
%
%   noprefix    -> noname                 the number or title alone, with no
%                                         name before it ("2", not
%                                         "Section 2")
%   S           -> S                      the first name capitalized and
%                                         unabbreviated, for a reference that
%                                         opens a sentence
%   capital     -> cap=true,abbrev=false  the names capitalized ("Lemma 1")
%   lowercase   -> cap=false,abbrev=false the names in lowercase ("lemma 1")
%   abbreviated -> abbrev=true            the abbreviated names ("Lem. 1")
%                                         where the reference type declares
%                                         them, the full names otherwise
%   anything else -> dropped              the default wording; a variant of
%                                         your own naming has no counterpart
%                                         in zref-clever, which would stop on
%                                         a key it does not know
%
% Redefine \flmRef to set references differently -- another package, another
% wording, your own variants, no link. Your definition sees the options as the
% document wrote them, in FLM's vocabulary, and maps them as it likes.
%
% The translation uses expl3, LaTeX's programming layer, which is part of the
% format and needs no package; its names read only between \ExplSyntaxOn and
% \ExplSyntaxOff. \keyval_parse:NNn reads the option list as LaTeX reads any
% key-value list: it hands a key written alone to the first macro below and a
% key with its value to the second, which drops the value. \str_case:nnF looks
% the key up in the table and appends what \zcref is to be given, or nothing.
% The macros with @ in their names are internals, not a customization point:
% redefine \flmRef, not its parts.
\ExplSyntaxOn
\def\flm@ref@add#1{\tl_put_right:Nn \flm@ref@options {#1,}}
\def\flm@ref@key#1{\str_case:nnF {#1}
  {{noprefix}{\flm@ref@add{noname}}
   {S}{\flm@ref@add{S}}
   {capital}{\flm@ref@add{cap=true,abbrev=false}}
   {lowercase}{\flm@ref@add{cap=false,abbrev=false}}
   {abbreviated}{\flm@ref@add{abbrev=true}}}
  {}}
\def\flm@ref@pair#1#2{\flm@ref@key{#1}}
\def\flm@ref@zcref#1#2{\zcref[#1]{#2}}
\ProvideDocumentCommand\flmRef{O{} m}{%
  \def\flm@ref@options{}%
  \keyval_parse:NNn \flm@ref@key \flm@ref@pair {#1}%
  \exp_args:No \flm@ref@zcref {\flm@ref@options} {#2}}
\ExplSyntaxOff
\endinput
